\chapter{从四足到双足点足：欠驱倒立摆、落脚平衡与硬约束 OCS2}
\label{ch:legged-biped}

% \cnum（字体无关带圈数字）已上提到 common/preamble.tex，全书统一，勿在此重定义。

\begin{practice}[前置自测]
开始之前，先试答下列问题。这些问题都是“四足栈移植到欠驱点足双足”时真实踩过的坑；若已能流畅作答，可只速读本章表格与速查卡；若卡住，正是本章要为你解决的。
\begin{enumerate}
  \item 四足小跑时两脚着地、绕支撑线欠驱（\cref{ch:quad_mpc}），但机器人靠 WBC 一个\emph{纯前馈}的机身（base）任务就能稳稳站住。把同一套 WBC 代码、同一个“base 反馈增益 $=0$”的配置原样搬到\emph{点足双足}上，它却会在 $1\,\mathrm{s}$ 内倒下。为什么四足能、点足不能？两者的支撑“多边形”差在哪？
  \item 一个直立的扫帚倒下去需要约 $1\,\mathrm{s}$。若一台点足机器人在 $50\,\mathrm{ms}$ 内俯仰角从 $0.004$ 冲到 $-0.25\,\mathrm{rad}$（比扫帚自然倒下快约 $17$ 倍），这说明问题出在“规划层不会走”还是“执行层在主动把它推倒”？你能用什么量化基线证明这个判断？
  \item 线性倒立摆（LIP）动力学是 $\ddot{x}=\omega^2 x$（注意是\emph{正号}）。这个正号意味着什么？能否找到一个随质心位置与速度变化的量 $\xi$，使它满足一阶线性方程 $\dot\xi=\omega\xi$，从而把“整个不稳定性”浓缩成一维？
  \item 点足只有一个点接触，踝部不能产生力矩去调压心（CoP）。那么当质心已经冲出脚掌附近、$\xi$ 以 $e^{\omega t}$ 发散时，唯一能“重置”它的物理手段是什么？为什么纯 PD 反馈救不回来、而\emph{迈一步}可以？
  \item 横向（左右方向）的平均速度明明是 $0$，为什么双足却必须以一个\emph{非零步宽} $W$ 左右交替地迈步、形成一个周期为 $2$ 的极限环？如果不主动给横向落脚反馈，会发生什么？
  \item OCS2 这类非线性 MPC 框架里，你给摩擦锥注册了一条\emph{硬不等式} $f^z\ge 50\,\mathrm{N}$，求解器也“收敛”了。这是否等于闭环里这条约束真被满足了？你该查哪个标量才能确认？
\end{enumerate}
\end{practice}

\section*{本章目标}
学完本章，你应当能够：
\begin{enumerate}
  \item \textbf{说清}点足双足为何\emph{没有静止站立的稳态}——支撑多边形退化成一个点（零面积），它是一个\emph{主动不稳的倒立摆}，“站立”在物理上只能是“踏步”（\cref{sec:lb-pointfoot}）；
  \item \textbf{推导}线性倒立摆（LIP）的不稳定动力学 $\ddot x=\omega^2 x$、分离出发散模态定义\emph{发散分量}（DCM）$\xi=x+\dot x/\omega$，并证明 $\dot\xi=\omega\xi$、$\xi(t)=\xi(0)e^{\omega t}$，从而把\emph{捕获点}（capture point）确立为“一步停稳”的落脚点（\cref{sec:lb-lip}）；
  \item \textbf{把}四足 gait 章里那个工程化的落足速度补偿项 $\sqrt{z_0/g}\,(\mathbf{v}-\mathbf{v}_{des})$（\cref{eq:qg-raibert-vel}）\textbf{认出来}——它正是捕获点反馈项，并据此补上它的 DCM 理论解释（\cref{sec:lb-lip}）；
  \item \textbf{掌握}横向（frontal plane）平衡的 ALIP 交替步宽闭式律，理解“横向平均速度为零却须以非零步宽 period-2 极限环行走”的根，并把它定位在 ALIP/H-LIP/DCM 的落脚律文献谱系里（\cref{sec:lb-alip}）；
  \item \textbf{建立} OCS2 的“OCP 五块拼图 + ReferenceManager”心智模型，弄清“为了并行 $\Rightarrow$ OCP 必须可拷贝 $\Rightarrow$ 不能直接改 problem $\Rightarrow$ 需 ReferenceManager/SynchronizedModule”这条设计因果链，并照清单把一条\emph{硬不等式}约束接入 OCP（\cref{sec:lb-ocs2}）；
  \item \textbf{辨清}“注册硬约束 $\ne$ 闭环满足”——理解 \texttt{sqp.g\_max} 收敛为何不含不等式残差、唯一证据是 \texttt{inequalityConstraintsSSE}，并看穿软圆锥的两个静默陷阱（\texttt{mu} 是惩罚权重不是摩擦系数、默认软锥对负法向力“免费”）（\cref{sec:lb-ocs2}）；
  \item \textbf{把}点足 WBC 的机身任务从“纯前馈”\textbf{拆分}为“base-XY 线加速度 + base 高度运动 + base 角运动”三块，写出角任务的正确形式（世界系旋转矩阵误差 \texttt{rotationErrorInWorld} + 世界系 $\boldsymbol\omega$ + base 角 Jacobian），并解释旧前馈写法的两块根因（切空间错配的 frame bug + 有限差分噪声放大 $\sim 476\times$）（\cref{sec:lb-wbc}）；
  \item \textbf{落地}运行时步态注入的 $\sim 1$ 个时域延迟问题（改 \texttt{defaultModeSequenceTemplate} 才能“从一开始就走”）、以及 \texttt{ModeSequenceTemplate} 平铺成交替步态的正确写法与“切换时刻数必比模式数多一个”的铁律（\cref{sec:lb-runtime}）。
\end{enumerate}

\subsection*{本章知识导航}
本章是一条“\emph{把已经跑通的四足栈（\cref{ch:quad_arch,ch:quad_gait,ch:quad_mpc,ch:quad_wbc}）移植到一台没有静稳的欠驱点足双足上}”的主线。它的灵魂不是“又一个新机器人”，而是一个反复出现的判断：\textbf{四足靠支撑多边形“接住”了很多本该由反馈完成的工作，点足把这层缓冲撤掉后，那些被默认关掉、被纯前馈化的环节就会逐一暴露}。结构如下：

\begin{center}
\begin{forest} navtreeLR
[{双足点足运动控制\\（四足栈$\to$欠驱点足）}
  [{点足为何特别\\支撑多边形$\to$一点}
    [{无静稳=倒立摆}] ]
  [{LIP/DCM/capture\\$\ddot x=\omega^2x$,\,$\dot\xi=\omega\xi$}
    [{认出 Raibert 反馈项}] ]
  [{ALIP 横向闭式律\\非零步宽 period-2}
    [{H-LIP/DCM 文献谱}] ]
  [{OCS2 硬约束接入\\五块拼图+RefManager}
    [{注册$\ne$满足}] ]
  [{点足 WBC 拆分\\base 角/高/XY 三块}
    [{frame bug+差分噪声}] ]
  [{运行时步态注入\\改默认序列}
    [{平铺陷阱}] ]
]
\end{forest}
\end{center}

\noindent\textbf{推荐路径}：\cref{sec:lb-pointfoot} 用“支撑多边形退化成一个点”一句立论，是全章的物理起点，任何读者都该读；\cref{sec:lb-lip,sec:lb-alip} 是“为什么迈步能稳、横向为什么必须交替迈步”的公式底座（LIP/DCM/capture/ALIP），它给四足 gait 章那个工程公式补上理论根，是理解后面一切工程决策的钥匙；\cref{sec:lb-ocs2,sec:lb-wbc} 是把上面这套移植进真实 OCS2+WBC 栈时\emph{真正会踩的两个工程深坑}（硬约束接入、机身任务拆分），写代码、排错前必须吃透；\cref{sec:lb-runtime} 是“怎么让它从第一拍就走起来”的运行时细节。

\subsection*{前置知识桥接}
本章把 P7 四足部与几何部的工具直接用作地基，凡前面已更深处讲过的，本章只在用到处 交叉引用 回引、不复述：
\begin{itemize}
  \item \cref{ch:quad_arch}（导论与控制架构）给出本章一切立论的总框架：浮动基\emph{欠驱}（\cref{ins:qa-floating-base}）、地面反力是\emph{唯一}可控外力（\cref{ins:qa-grf-only}）、多模型须切换（\cref{ins:qa-multimodel}）。点足平衡正是把这三条推到极限的产物；
  \item \cref{ch:quad_gait}（步态与落足）§\ref{sec:qg-footstep} 已写出 Raibert 三项分解，其中的速度补偿项 $\sqrt{z_0/g}\,(\mathbf{v}-\mathbf{v}_{des})$（\cref{eq:qg-raibert-vel}）与“倒立摆中性点在髋正下方”（\cref{ins:qg-neutral-why}）正是本章 DCM/捕获点理论的工程影子，\cref{sec:lb-lip} 给它补理论；步态调度器 ModeSchedule（\cref{sec:qg-scheduler}）是 \cref{sec:lb-runtime} 平铺陷阱的对象；
  \item \cref{ch:quad_mpc}（单刚体 MPC）§\ref{sec:srbd} 已把 LIPM（Kajita）/ 质心动力学（Orin）/ 单刚体三模型并置对照、并给出凸化与摩擦金字塔（\cref{sec:convexify}）；本章 LIP/DCM 接其 LIPM，硬约束接其 QP 化（\cref{sec:qp,sec:state-space}），软/硬摩擦约束接其凸化；
  \item \cref{ch:quad_wbc}（全身控制 WBC）§\ref{sec:qw-floating} 已建立浮动基 WBC 的子任务体系、零空间投影与松弛 QP（\cref{sec:qw-relax}）；本章 \cref{sec:lb-wbc} 是这套子任务在“点足无静稳”下的\emph{特化}——尤其机身加速度任务（\cref{tab:qw-subtasks} 中那一项）为何对点足必须带反馈；
  \item \cref{ch:lie}（李群与李代数）给出 $\dot{\mathbf{R}}=\boldsymbol{\omega}^\wedge\mathbf{R}$ 与旋转误差的几何，本章机身角任务的 \texttt{rotationErrorInWorld} 与“为何不能直接在欧拉角空间相加”都建在这条微分方程上（亦见 \cref{sec:qw-practice} 的欧拉角误差变基）。
\end{itemize}

\subsection*{如果跳过本章会怎样}
\begin{itemize}
  \item 你会以为“四足能站 $\Rightarrow$ 同一套 WBC 移到双足也能站”，从而把整条调试战线压在 OCS2 上游（落脚/步时/力可行性），而真正的根因——执行层一个被悄悄置零的机身反馈增益——可能让你花上百个调试回合都摸不到（\cref{sec:lb-wbc} 即这条血泪史的方法论提炼，完整复盘见 \cref{ch:legged-debug}）；
  \item 你会把四足 gait 章那个落足速度补偿项当成一个“拍脑袋的工程公式”照抄，而不知道它就是\emph{捕获点}、背后是 $\dot\xi=\omega\xi$ 这条把全部不稳定性浓缩成一维的方程——于是不会调它、不会在它失效时（横向）补上 ALIP 闭式律。
\end{itemize}

\begin{table}[H]\centering\small
\caption{本章阅读方式建议}
\begin{tabular}{lll}
\toprule
阅读方式 & 时间 & 适合谁 \\
\midrule
精读（含全部推导与练习） & 4--5 小时 & 首次系统学习、要做双足/点足移植 \\
速读（跳过 ALIP 闭式推导与 WBC 细节） & 1.5 小时 & 有四足 MPC/WBC 基础、想对齐点足差异 \\
速查（只看小结的符号表 + 公式速查表） & 15 分钟 & 写控制器代码时回来查 LIP/DCM/capture/ALIP \\
\bottomrule
\end{tabular}
\end{table}

\begin{note}
\textbf{本章记号与约定.}\;
沿用 \cref{ch:quad_mpc} 的记号：旋转矩阵 $\mathbf{R}\in\SO(3)$、ZYX 欧拉角 $\boldsymbol{\Theta}=[\phi,\theta,\psi]^\top$（roll/pitch/yaw）、反对称算子 $(\cdot)^\wedge$ 与其逆 $(\cdot)^\vee$（\cref{ch:lie}）。
\textbf{倒立摆量}：质心高度记 $h$（或恒值 $z_0$）、质心水平位置 $x$（矢状面）/ $y$（横向）、LIP 自然频率 $\omega=\sqrt{g/h}$；横向 LIP 频率沿用源记号写作 $\ell$。
\textbf{发散分量（DCM）/ 捕获点}记 $\xi$（同一符号兼用于“瞬时捕获点”）；支撑足（落脚点）水平位置记 $p$（矢状）/ $\mathbf{p}_{\mathrm{stance}}$（向量）。
\textbf{步态量}：步时（一步时长）记 $T$、步宽记 $W$、横向落脚量记 $L_x^{des}$（沿用源记号，下标 $x$ 指“横向位移量”而非世界系 $x$ 轴）。
\textbf{维度}：点足双足的质心状态维 \texttt{stateDim}$=18$、输入维 \texttt{inputDim}$=12$、WBC 决策维 $24$、接触数 \texttt{numContacts}$=2$（对照四足 $24/24/42/4$，\cref{ch:quad_mpc}）。
本章数值取自一台代号 Braver 的点足双足平台：$h\approx 0.427\,\mathrm{m}$、$\omega\approx 4.8\,\mathrm{rad/s}$、总质量 $\approx 9.8\,\mathrm{kg}$、横向两脚相距 $\pm 0.1475\,\mathrm{m}$ 但前后位于\emph{同一} $x$ 线（$x\approx 0.0235\,\mathrm{m}$）。
\end{note}

% ════════════════════════════════════════════════════════════════
\section{为什么点足“特别”：从静稳多边形到主动不稳倒立摆}
\label{sec:lb-pointfoot}

\paragraph{动机：把同一套四足栈搬到点足，第一性差异是什么？}
本章的整个工程语境是“移植”：一套基于 OCS2 质心 NMPC + WBC 的四足控制栈（\cref{ch:quad_arch,ch:quad_mpc,ch:quad_wbc}）已经能让四足稳稳行走，现在要把它\emph{原样}用到一台只有两条腿、每条腿末端是一个\emph{点}的双足上。移植腿式机器人有一长串通用坑（关节顺序、四元数约定、URDF 决定 Pinocchio 索引、维度从 \texttt{CentroidalModelInfo} 读取而非硬编码等，\cref{ch:legged-debug} 系统整理），但本章只聚焦那个\emph{由“点足”这一几何事实直接逼出来、四足身上根本不会暴露}的核心差异。它一句话可说清：\textbf{四足靠一个有面积的支撑多边形“站得住”，点足的支撑多边形退化成一个点（零面积），于是它根本没有静止站立的稳态——它是一个主动不稳的倒立摆，“站立”只能是“踏步”。}

\paragraph{维度对照：先看“账面”上点足比四足少了什么。}
表面上点足只是四足的“缩水版”：\cref{tab:lb-dims} 把两者的质心状态/输入维度并排。但这张表的真正信息不在数字大小，而在最后一行——\textbf{接触数从 $4$ 掉到 $2$}，而且这 $2$ 个接触点的几何排布（下文）让“支撑多边形”这个四足赖以稳定的概念彻底失效。

\begin{table}[H]\centering\small
\caption{四足与点足双足的质心模型维度对照（均按 \cref{ch:quad_mpc} 的质心状态约定，维度从 \texttt{CentroidalModelInfo} 读取）}
\begin{tabular}{lcc}
\toprule
量 & 四足（如 Go1） & 点足双足（Braver） \\
\midrule
驱动自由度（每腿关节数$\times$腿数） & $3\times4=12$ & $3\times2=6$ \\
质心状态维 \texttt{stateDim}$=6+6+$驱动DoF & $24$ & $18$ \\
输入维 \texttt{inputDim}$=3\cdot$\texttt{numContacts}$+$驱动DoF & $24$ & $12$ \\
WBC 决策维 $=$ 广义坐标$+3\cdot$\texttt{numContacts}$+$驱动DoF & $42$ & $24$ \\
\textbf{接触数 \texttt{numContacts}} & $\mathbf{4}$ & $\mathbf{2}$ \\
支撑多边形 & 四脚张成的\emph{二维区域} & 退化为\emph{一个点/一条短线} \\
\bottomrule
\end{tabular}
\label{tab:lb-dims}
\end{table}

\paragraph{核心立论：两脚同一 $x$ 线 $\Rightarrow$ 俯仰方向零力臂 $\Rightarrow$ 无静稳。}
这台点足双足的两脚在横向相距 $\pm 0.1475\,\mathrm{m}$，但\emph{前后位于同一条 $x$ 线上}（两个接触点的 $x$ 坐标都约为 $0.0235\,\mathrm{m}$）。把这两个点接触连起来，得到的不是一个有面积的多边形，而是一条\emph{横向的短线段}。于是：

\begin{itemize}
  \item \textbf{横向（roll/$y$）}：两脚相距 $0.295\,\mathrm{m}$，对横向倾倒\emph{有}力臂——但这条线无限薄，能提供的横向恢复力矩完全来自两脚法向力之差，而点接触又不能产生踝力矩，所以横向稳定也极脆；
  \item \textbf{矢状（pitch/$x$）}：两脚\emph{同一} $x$ 坐标 $\Rightarrow$ 对俯仰方向\emph{完全没有支撑力臂}。无论两脚用多大的法向力，都\emph{不可能}产生一个绕横轴（$y$ 轴）的恢复力矩去阻止机身前扑或后仰。
\end{itemize}

第二条是致命的。回忆 \cref{ch:quad_mpc} 里四足对角步态“绕支撑线欠驱”——那时机器人\emph{也}是两脚着地、绕支撑线方向没有力臂。但四足的支撑线半个周期就切换到另一对脚，绕新支撑线\emph{又}欠驱，靠的是 MPC 在欠驱来临前预先准备（\cref{ins:wbc-no-predict}）。点足的处境更极端：它\emph{始终}只有两个同 $x$ 的接触点，俯仰方向\emph{永远}没有静止支撑力臂——这不是“间歇欠驱”，而是“\emph{结构性}地没有静稳平衡点”。

\begin{insight}[点足无静稳：从“浮动基欠驱”推到极限]\label{ins:lb-no-static}
\cref{ins:qa-floating-base} 说浮动基系统是欠驱的——机身的 $6$ 个自由度没有任何关节电机直接驱动，全靠地面反力（\cref{ins:qa-grf-only} 进一步指出地面反力是\emph{唯一}可控外力）。四足也欠驱，但它被四脚张成的二维支撑多边形“接住”了：CoM 只要落在多边形内，重力线就被支撑域包住，姿态反馈即使做得很弱、很抖，机器人也不会真的倒——\emph{多边形替反馈兜了底}。点足把这个多边形压成一个点（俯仰方向甚至是零维），就等于\textbf{撤掉了这层兜底}。此时机身姿态的唯一抓手是穿过那个固定点接触的地面反力，受摩擦锥限、且对俯仰无力臂；只要不主动、持续地反馈纠正，机身就会像扫帚一样倒下去。\textbf{“点足没有静止站立稳态、站立必须是踏步”是本章一切工程决策的物理前提。}
\end{insight}

\begin{figure}[ht]
\centering
\begin{tikzpicture}[scale=1.0,>=Stealth]
  % ===== left: quadruped support polygon (area) =====
  \begin{scope}[shift={(-3.6,0)}]
    \node[main!60!black] at (0,1.9) {\small 四足：支撑多边形\emph{有面积}};
    % four feet
    \coordinate (a) at (-1.1,0.7);
    \coordinate (b) at (1.1,0.5);
    \coordinate (c) at (1.0,-0.9);
    \coordinate (d) at (-1.0,-0.7);
    \fill[main!12,draw=main!60!black,thick] (a)--(b)--(c)--(d)--cycle;
    \foreach \p in {a,b,c,d} \fill (\p) circle (2.2pt);
    % CoM projection inside
    \fill[red!70!black] (-0.05,-0.05) circle (1.8pt) node[above right,red!70!black]{\scriptsize CoM 投影};
    \node[main!60!black,align=center] at (0,-1.6) {\scriptsize CoM 落在多边形内\\\scriptsize $\Rightarrow$ 有静稳平衡};
  \end{scope}
  % ===== right: point-foot biped degenerate =====
  \begin{scope}[shift={(3.6,0)}]
    \node[second!50!black] at (0,1.9) {\small 点足：支撑“多边形”退化成一条线};
    % two feet on same x-line, separated only laterally
    \coordinate (l) at (0,0.8);
    \coordinate (r) at (0,-0.8);
    \draw[second,very thick,dash dot] (l)--(r);
    \fill (l) circle (2.2pt) node[left]{\scriptsize 左脚};
    \fill (r) circle (2.2pt) node[left]{\scriptsize 右脚};
    % CoM
    \fill[red!70!black] (0.0,0.0) circle (1.8pt);
    % lateral arm (has lever)
    \draw[teal!60!black,<->] (0.35,0.8)--(0.35,-0.8) node[midway,right,teal!60!black]{\scriptsize 横向有力臂};
    % sagittal: zero arm — arrow into the line
    \draw[third,-{Stealth[length=2.5mm]},thick] (-1.3,0)--(-0.15,0);
    \node[third,align=center] at (-1.5,-0.6) {\scriptsize 俯仰方向\\\scriptsize 零力臂$\Rightarrow$必前扑};
    \node[second!50!black,align=center] at (0,-1.6) {\scriptsize 同 $x$ 线\,$\Rightarrow$\,无静稳\\\scriptsize $\Rightarrow$“站立”=“踏步”};
  \end{scope}
\end{tikzpicture}
\caption{四足 vs 点足的支撑域对比。左：四脚张成有面积的支撑多边形，CoM 投影落在域内即有静稳平衡，多边形替姿态反馈“兜底”。右：点足双足两脚位于\emph{同一} $x$ 线、仅横向分开，支撑域退化成一条无限薄的横线——横向尚有力臂（但极脆），\emph{矢状（俯仰）方向零力臂}，故没有静止站立稳态，机身必前扑：稳定只能靠迈步（\cref{sec:lb-lip,sec:lb-alip}）。}
\label{fig:lb-support}
\end{figure}

\begin{pitfall}[误区：把四足“base 反馈可选”的直觉搬到点足]\label{pit:lb-quad-intuition}
四足上有一个极具迷惑性的经验事实：WBC 的机身（base）加速度任务即使\emph{只做纯前馈}（用质心动量算出期望机身加速度，但\emph{不}加任何位姿/速度反馈，等价于反馈增益 $=0$），机器人也能稳稳站住、甚至小跑。于是一个非常自然却致命的推论是：“机身反馈是可选的、锦上添花的”。这条直觉在四足上成立\emph{仅仅因为}支撑多边形替它兜了底（\cref{ins:lb-no-static}）。把这个“反馈增益 $=0$”的配置原样移植到点足，等于把点足\emph{唯一}的姿态反馈通道掐断——倒立摆失去主动镇定，必然发散。这正是本章 \cref{sec:lb-wbc} 要彻底纠正的根因，也是“移植腿式机器人时必须质疑一切被默认置零的反馈增益”这条铁律的来源（完整复盘见 \cref{ch:legged-debug}）。\textbf{记住：在点足上，机身反馈不是可选项，是续命项。}
\end{pitfall}

\begin{practice}
\begin{enumerate}
  \item 用 \cref{ins:qa-grf-only}（地面反力是唯一外力）说明：为什么点足在俯仰方向“零力臂”意味着\emph{无论}两脚法向力如何分配，都产生不了俯仰恢复力矩？（提示：力臂为零的力对该轴的力矩恒为零。）
  \item 对照 \cref{fig:underactuated}（四足绕支撑线欠驱）与本章 \cref{fig:lb-support}：四足的欠驱是“间歇”的（半周期切换），点足的“无静稳”是“结构性”的。这两种性质对“控制器是否必须始终在踏步”的要求有何不同？
  \item （概念）若把这台点足的两脚改成有一定前后间距（不再同 $x$），俯仰方向就有了力臂、能像“平足”一样静站。这说明“点足无静稳”究竟是几何属性还是动力学属性？它能否通过\emph{换脚型}而非\emph{换控制律}解决？
\end{enumerate}
\end{practice}

% ════════════════════════════════════════════════════════════════
\section{线性倒立摆、发散分量与捕获点}
\label{sec:lb-lip}

\paragraph{动机：既然没有静稳，就需要一个“迈步能稳”的最小模型。}
\cref{sec:lb-pointfoot} 确立了点足是倒立摆、平衡只能靠迈步。但“迈到哪里”才能稳？要回答这个，需要一个比单刚体 MPC（\cref{ch:quad_mpc}）更精简、能给出\emph{闭式}落脚律的模型。这个模型就是\textbf{线性倒立摆}（Linear Inverted Pendulum, LIP），它正是 \cref{sec:srbd} 三模型谱系里最左、最简的那个（Kajita\cite{kajita2001lipm}），\cref{ch:quad_mpc} 在那里只把它作为“质点无转动”的谱系起点一笔带过；本节把它\emph{推到底}，因为它对点足双足不是简化的简化，而是平衡问题的\emph{核心数学结构}。

\subsection{LIP 动力学：那个致命的正号}
\label{subsec:lb-lip-dyn}

把点足双足简化成“一个质点（质心，质量 $m$）+ 一根无质量伸缩腿”，腿的下端固连在支撑点 $p$（落脚点的水平位置）。约束质心高度恒为 $h$（这是 LIP 的定义性假设，使动力学线性）。在矢状面内，质心相对支撑点的水平位移为 $(x-p)$，沿腿方向的支撑力与重力共同作用，水平方向的运动方程为

\begin{equation}
  \ddot{x}=\frac{g}{h}\,(x-p)=\omega^2\,(x-p),
  \qquad
  \omega:=\sqrt{\frac{g}{h}}.
  \label{eq:lb-lip}
\end{equation}

把支撑点取作原点（$p=0$）即得最干净的形态 $\ddot x=\omega^2 x$。\cref{ch:quad_mpc} 里我们反复处理的刚体动力学都是“受控就能镇定”的；而这里出现了一个\emph{决定一切}的细节——\textbf{是正号 $\omega^2 x$，不是负号}。负号（$\ddot x=-\omega^2 x$）是简谐振子，稳定、来回振荡；正号意味着特征根是\emph{实}的 $\pm\omega$，系统有一个\emph{发散}模态：质心离支撑点越远，被推得越远。这就是“倒立摆”三个字的数学本体——一个天然不稳定的二阶系统。对 Braver，$h\approx 0.427\,\mathrm{m}$、$g=9.81\,\mathrm{m/s^2}$，故 $\omega=\sqrt{9.81/0.427}\approx 4.8\,\mathrm{rad/s}$。

\begin{insight}[正号 $\Rightarrow$ 实特征根 $\pm\omega$ $\Rightarrow$ 必有发散模态]\label{ins:lb-positive-sign}
线性系统 $\ddot x=\omega^2 x$ 的通解是 $x(t)=A\,e^{\omega t}+B\,e^{-\omega t}$：一个发散模态 $e^{\omega t}$、一个收敛模态 $e^{-\omega t}$。任何初始扰动只要在 $e^{\omega t}$ 方向有分量，就会指数级放大。\textbf{倒立摆的全部不稳定性都装在 $e^{\omega t}$ 这一个模态里}——这给了我们一个极强的简化希望：如果能找到一个\emph{只}沿发散方向的标量 $\xi$，把整个二维相平面 $(x,\dot x)$ 的不稳定性压缩成一维 $\dot\xi=\omega\xi$，控制问题就从“镇定一个二阶不稳定系统”变成“把一个一维标量摁住”。这正是下一小节 DCM 的动机。
\end{insight}

\subsection{发散分量（DCM）：把不稳定性浓缩成一维}
\label{subsec:lb-dcm}

为分离 \cref{eq:lb-lip}（取 $p=0$）的发散模态，定义\textbf{发散分量}（Divergent Component of Motion, DCM），在双足/人形文献里也叫\emph{瞬时捕获点}，记

\begin{equation}
  \boxed{\;\xi := x+\frac{\dot x}{\omega}\;}
  \label{eq:lb-dcm-def}
\end{equation}

这个组合不是凑出来的：$\xi$ 恰好就是相平面里沿 $e^{\omega t}$ 特征方向的坐标。对它求导，并代入 LIP 动力学 $\ddot x=\omega^2 x$：

\begin{equation}
  \dot\xi=\dot x+\frac{\ddot x}{\omega}
  =\dot x+\frac{\omega^2 x}{\omega}
  =\dot x+\omega x
  =\omega\Big(x+\frac{\dot x}{\omega}\Big)
  =\omega\,\xi.
  \label{eq:lb-dcm-dyn}
\end{equation}

于是得到一阶线性方程 $\dot\xi=\omega\xi$，解为

\begin{equation}
  \xi(t)=\xi(0)\,e^{\omega t}.
  \label{eq:lb-dcm-sol}
\end{equation}

\textbf{这一步是整套双足平衡理论的枢纽}（DCM 的三维推广由 Englsberger 等系统化\cite{englsberger2015dcm}）。它说：原来纠缠在一起的二阶不稳定动力学，被坐标变换 \cref{eq:lb-dcm-def} 干净地拆成了“一个一维发散的 $\xi$（\cref{eq:lb-dcm-sol}）+ 一个被它带着走的剩余收敛分量”。我们\emph{不必}去镇定整个 $(x,\dot x)$ 平面——只要控制住 $\xi$，就控制住了全部不稳定性。剩下那个分量（常记 $x-\xi$，即所谓“稳定分量”CoM 的收敛部分）自己会衰减，无须操心。

\begin{note}[DCM 与质心速度的关系、以及它为何叫“瞬时捕获点”]\label{note:lb-dcm-name}
把 \cref{eq:lb-dcm-def} 读成 $\xi=x+\dot x/\omega$：它是“当前质心位置”加上“一个正比于质心速度的前瞻量 $\dot x/\omega$”。直觉上，$\xi$ 是“若此刻不再施加任何控制、质心将渐近停在的那个水平点”——质心冲得越快（$\dot x$ 越大），这个点越靠前。这正是“捕获点”一词的来源：\emph{把脚迈到这个点上，机器人就能一步把动能耗尽、停下来}。因为支撑点在 $\xi$ 处时，$\xi$ 相对支撑点为零，发散模态被置零，倒立摆“被捕获”。下一小节把这点说严格。
\end{note}

\subsection{捕获点：一步停稳的落脚点}
\label{subsec:lb-capture}

现在把落脚点 $p$ 放回来（不再设 $p=0$）。把 \cref{eq:lb-lip} 完整版 $\ddot x=\omega^2(x-p)$ 重做一遍 DCM 推导：定义 $\xi=x+\dot x/\omega$（与支撑点无关，仍是 \cref{eq:lb-dcm-def}），则

\begin{equation}
  \dot\xi=\dot x+\frac{\ddot x}{\omega}
  =\dot x+\omega(x-p)
  =\omega\Big(x+\frac{\dot x}{\omega}-p\Big)
  =\omega\,(\xi-p).
  \label{eq:lb-dcm-stance}
\end{equation}

这是带支撑点的 DCM 动力学。它立刻给出\textbf{捕获点}（capture point）的定义：令发散停止，即令 $\dot\xi=0$，需要

\begin{equation}
  \boxed{\;\xi-p=0\quad\Longrightarrow\quad p_{\mathrm{capture}}=\xi=x+\frac{\dot x}{\omega}\;}
  \label{eq:lb-capture}
\end{equation}

即\emph{把支撑点迈到当前 DCM 处}。一旦 $p=\xi$，\cref{eq:lb-dcm-stance} 给出 $\dot\xi=0$，DCM 冻结、发散模态被中和，机器人在这一步内停稳。这就是“捕获”的精确含义。\cref{fig:lb-dcm} 把这套关系画出来。

\begin{figure}[ht]
\centering
\begin{tikzpicture}[scale=1.05,>=Stealth]
  % ground
  \draw[gray] (-0.5,0)--(6.2,0);
  \foreach \x in {-0.2,0.2,...,6.0} \draw[gray] (\x,0)--(\x-0.13,-0.12);
  % stance point p
  \fill (1.2,0) circle (2.4pt) node[below]{\scriptsize 支撑点 $p$};
  % CoM
  \coordinate (com) at (2.2,2.0);
  \fill[main!70!black] (com) circle (2.6pt) node[above]{\scriptsize 质心 $x$ (高 $h$)};
  % massless leg
  \draw[gray!60!black,very thick] (1.2,0)--(com);
  % CoM velocity
  \draw[red!70!black,-{Stealth[length=2.5mm]},thick] (com)--++(1.4,0) node[right,red!70!black]{\scriptsize $\dot x$};
  % DCM xi = x + xdot/omega  (project to ground for picture)
  \coordinate (xi) at (3.9,0);
  \draw[teal!60!black,dashed] (com)--(xi);
  \fill[teal!60!black] (xi) circle (2.4pt) node[below]{\scriptsize $\xi=x+\dot x/\omega$ (=捕获点)};
  % brace x to xi
  \draw[decorate,decoration={brace,amplitude=4pt},teal!60!black] (2.2,-0.55)--(3.9,-0.55) node[midway,below,teal!60!black]{\scriptsize $\dot x/\omega$ 前瞻量};
  % annotation
  \node[align=center,second!50!black] at (3.0,1.0) {\scriptsize 把脚迈到 $\xi$\\\scriptsize $\Rightarrow\dot\xi=0$ 一步停稳};
\end{tikzpicture}
\caption{捕获点几何。质心 $x$ 以速度 $\dot x$ 冲出，DCM $\xi=x+\dot x/\omega$ 是“位置 + 正比于速度的前瞻量”所指的地面点。把支撑点 $p$ 迈到 $\xi$ 上（\cref{eq:lb-capture}），则 \cref{eq:lb-dcm-stance} 给出 $\dot\xi=0$，发散模态被中和、机器人一步停稳；若想继续走，则把脚落在比 $\xi$ 稍\emph{偏后}处，留一点残余 DCM 驱动下一步。}
\label{fig:lb-dcm}
\end{figure}

\begin{insight}[为何纯 PD 镇定 DCM 而不迈步行不通：唯一抓手受困于固定点接触]\label{ins:lb-must-step}
读者可能会问：既然不稳定性已浓缩成一维 $\dot\xi=\omega(\xi-p)$，为什么不直接对 $\xi$ 做反馈、用地面反力把它拉回来，而非要\emph{迈步}？关键在于\emph{抓手受限}。能影响 $\xi$ 的物理量只有地面反力，而对点足，这个力\textbf{必须穿过那个固定的点接触}——它受摩擦锥约束（切向力不能超过 $\mu$ 倍法向力，\cref{sec:convexify}），而且因为是点接触、踝部无力矩，\emph{不能}像平足那样在脚掌内移动压心（CoP）来调节力矩。于是 $p$ 在一步之内基本是\emph{钉死}的。\cref{eq:lb-dcm-stance} 里 $p$ 固定时，$\xi$ 一旦跑出脚附近，$e^{\omega t}$（\cref{eq:lb-dcm-sol}）的发散速度很快超出地面反力在摩擦锥内能提供的修正能力——纯 PD 救不回来。\textbf{唯一能从根本上重置 $\xi$ 的手段，是移动支撑点本身，即迈步：让 $p$ 跳到新的 DCM 附近，重新令 $\xi-p$ 变小。}这就是“点足无静稳、横向平衡只能靠落脚”这句话的数学根（\cref{ins:lb-no-static} 的定量版）。
\end{insight}

\subsection{回看四足 gait 章：那个工程公式就是捕获点反馈}
\label{subsec:lb-raibert-link}

现在可以兑现本章开头埋下的承接钩子。\cref{ch:quad_gait} 在讲落足点规划时（\cref{sec:qg-footstep}）给出了 Raibert 式落足公式的三项分解，其中速度补偿项写作（\cref{eq:qg-raibert-vel}）

\begin{equation}
  \Delta\mathbf{p}_{vel}=\sqrt{\frac{z_0}{g}}\,(\mathbf{v}-\mathbf{v}_{des}),
  \label{eq:lb-raibert-recall}
\end{equation}

当时它是作为一个“鲁棒但非最优”的工程启发式给出的（\cref{ins:qg-raibert-nature}），并配了一条洞见说它是“捕获点刹车”（\cref{ins:qg-capture}）。现在我们能把这条洞见说\emph{严格}了：把 \cref{eq:lb-capture} 的捕获点 $p_{\mathrm{capture}}=x+\dot x/\omega$ 改写成“落脚点相对质心当前位置的\emph{偏移}”——令 $\mathbf{v}_{des}=0$（期望停下），偏移量 $p_{\mathrm{capture}}-x=\dot x/\omega=\dot x\,\sqrt{h/g}$。把 $h$ 记作 $z_0$、$\dot x$ 记作 $\mathbf{v}$，这恰好就是 \cref{eq:lb-raibert-recall} 里的 $\sqrt{z_0/g}\,\mathbf{v}$！更一般地，当期望以 $\mathbf{v}_{des}$ 巡航而非停下时，要落在“相对中性点对称”的位置，速度项就成了 $\sqrt{z_0/g}\,(\mathbf{v}-\mathbf{v}_{des})$。

\begin{insight}[同一个 $\sqrt{z_0/g}$：工程公式与 DCM 理论的会师]\label{ins:lb-raibert-is-capture}
\cref{ch:quad_gait} 那个看似“拍脑袋”的速度补偿项 $\sqrt{z_0/g}\,(\mathbf{v}-\mathbf{v}_{des})$，其系数 $\sqrt{z_0/g}=1/\omega$ 正是 LIP 自然频率的倒数，而整项正是\emph{捕获点相对中性点的偏移}（\cref{eq:lb-capture}）。换句话说：\cref{ch:quad_gait} 是从“让落点对称于倒立摆中性点”的几何直觉（\cref{ins:qg-neutral-why}：中性点在髋正下方、重力线过支点）\emph{反推}出这个公式的；本章则是从 LIP/DCM 动力学\emph{正推}出同一个公式。两条路殊途同归，这不是巧合——它印证了“用落足点控平衡”这件事的底层就是捕获点动力学（Raibert 1986\cite{raibert1986legged} 的落足启发式、Pratt 等\cite{pratt2006capture} 的捕获点框架是同一脉络）。理解这一点，你才知道 \cref{eq:lb-raibert-recall} 里每个符号从哪条物理方程来、$z_0$ 变了该怎么调、它在什么时候（横向、大扰动）会失效而需要升级（\cref{sec:lb-alip}）。
\end{insight}

\begin{remark}[本项目的 MPC 并不显式算捕获点，但隐式在做同一件事]\label{rem:lb-implicit-capture}
要澄清一个容易的误解：本章移植的 OCS2 质心 NMPC 栈\emph{并不}在代码里显式计算 \cref{eq:lb-capture} 这个公式，更不会把它当作硬目标喂给求解器。它做的是把落脚点的 XY 坐标作为\emph{优化变量}，在质心动力学 + 摩擦锥 + 步态时序的约束下，让 SQP 自己解出“为了跟踪期望速度且不打滑，脚该落在哪”。但在数学上，这个隐式优化所收敛到的落脚点，与捕获点动力学要求的是\emph{同一类}东西——求解器是在做“capture-point-like”的事，只不过把闭式律换成了带约束的数值优化（这也是 \cref{sec:lb-runtime} 里“放手让 MPC 自己拥有落点”最终能走通的原因）。把闭式捕获点理论握在手里的价值，不在于替代 MPC，而在于\emph{诊断}：当 MPC 解出的落脚明显偏离捕获点直觉时，你知道该去查哪里。
\end{remark}

\begin{practice}
\begin{enumerate}
  \item 由 \cref{eq:lb-lip} 完整版（含 $p$）重新推 \cref{eq:lb-dcm-stance}，确认 $\xi$ 的定义 \cref{eq:lb-dcm-def} 不随 $p$ 改变、且 $\dot\xi=\omega(\xi-p)$。再令 $p=\xi$ 验证 $\dot\xi=0$。
  \item 取 Braver 数值 $\omega=4.8\,\mathrm{rad/s}$。设质心当前在支撑点正上方（$x=0$）但有前向速度 $\dot x=0.3\,\mathrm{m/s}$，算出捕获点 $\xi$ 的位置（相对支撑点）。若想以 $\mathbf{v}_{des}=0.3\,\mathrm{m/s}$ 继续巡航，落脚应比 $\xi$ 偏前还是偏后？
  \item （承接）把 \cref{eq:lb-capture} 的偏移量 $\dot x/\omega$ 与 \cref{ch:quad_gait} 的 \cref{eq:lb-raibert-recall} 逐符号对齐，写清 $z_0\leftrightarrow h$、$\mathbf{v}\leftrightarrow\dot x$、系数 $\sqrt{z_0/g}\leftrightarrow 1/\omega$ 的对应。这道题的目的是让你\emph{亲手}确认两章用的是同一个公式（\cref{ins:lb-raibert-is-capture}）。
  \item （概念）\cref{ins:lb-must-step} 说纯 PD 镇定 $\xi$ 不迈步行不通，根因是地面反力受摩擦锥限且 $p$ 钉死。如果换成\emph{平足}（脚掌有面积、踝部能产生力矩、CoP 可在脚掌内移动），这个结论会松动吗？这解释了为什么 ZMP/CoP 框架是\emph{平足}双足的主流、而点足/欠驱双足更依赖落脚。
\end{enumerate}
\end{practice}

% ════════════════════════════════════════════════════════════════
\section{ALIP 横向落脚闭式律：为什么必须左右交替迈步}
\label{sec:lb-alip}

\paragraph{动机：矢状面解决了，横向（roll）还有一个独立的难题。}
\cref{sec:lb-lip} 用 LIP/DCM/捕获点把\emph{矢状面}（前后、pitch）的平衡讲透了。但横向（左右、roll）有一个矢状面没有的诡异之处：\textbf{横向的平均速度是 $0$}（机器人直线行走时不该往左或往右净移动），\emph{然而它却不能像“原地不动”那样横向静止}——双足必须以一个\emph{非零步宽} $W$ 左右交替地落脚。本节回答“为什么”，并给出横向落脚的\emph{闭式}律。

\subsection{横向 period-2 极限环：非零步宽从何而来}
\label{subsec:lb-lateral}

考虑横向的 LIP：质心横向坐标 $y$、横向支撑点（左脚或右脚的 $y$ 坐标）。横向动力学与 \cref{eq:lb-lip} 同形，只是频率记作 $\ell$（横向 LIP 频率，本台机器人 $\ell$ 与 $\omega$ 数值相近，因高度同为 $h$）。问题的关键在\emph{支撑点的交替}：左脚支撑时，支撑点在 $y=+W/2$ 一侧；右脚支撑时在 $y=-W/2$ 一侧。质心要直线前进、横向不净移，就必须在左右两个支撑点之间\emph{来回摆}——这是一个\textbf{周期为 $2$ 的极限环}（period-2 limit cycle）：第 $k$ 步落左、第 $k+1$ 步落右、第 $k+2$ 步又落左……横向状态每\emph{两}步才回到同一相位。

为什么不能步宽 $W=0$（两脚落在同一条中线上，像走钢丝）？因为 $W=0$ 时左右支撑点重合、横向支撑点不再交替，横向 LIP 退化成 \cref{eq:lb-lip} 那个纯发散系统，roll 模态没有任何东西去周期性地“接住”它——任何微小的横向扰动都会沿 $e^{\ell t}$ 发散倒向一侧。\textbf{非零步宽 $W$ 是横向稳定的结构性必需}：它让支撑点周期性地切换到质心倾倒的\emph{那一侧}去接住它，正如 \cref{ins:lb-must-step} 所说“移动支撑点重置 DCM”，只是横向的重置是\emph{周期、双侧交替}的。

\subsection{交替步宽闭式律}
\label{subsec:lb-alip-law}

要维持这个 period-2 极限环、且让横向 DCM 在每步首尾匹配（保证周期性而不发散），可以解出横向落脚量的\emph{闭式}表达式。在绕接触点角动量描述的 LIP（即 ALIP, Angular-momentum Linear Inverted Pendulum，Gong \& Grizzle\cite{gong2022alip}）框架下，维持步宽 $W$、步时 $T$ 的交替横向落脚量满足

\begin{equation}
  \boxed{\;L_x^{des}=\pm\frac{1}{2}\,mHW\cdot\frac{\ell\,\sinh(\ell T)}{1+\cosh(\ell T)}\;}
  \label{eq:lb-alip}
\end{equation}

其中 $\pm$ 随左右脚交替、$m$ 为质量、$H$ 为质心高度（即 $h$）、$W$ 为步宽、$T$ 为步时、$\ell$ 为横向 LIP 频率。式中 $\sinh,\cosh$ 的出现是 LIP 解 $e^{\pm\ell t}$ 的双曲函数组合（正号系统的自然产物，对照 \cref{ins:lb-positive-sign}）；$\frac{\ell\sinh(\ell T)}{1+\cosh(\ell T)}$ 这一因子把“一步时长 $T$ 内横向 DCM 应有的变化”折算成所需的落脚反馈量。直觉上：步时越长（$T$ 大）、步宽越大（$W$ 大），横向要“接住”的发散越多，所需落脚量越大。

\begin{insight}[横向 roll 模态没有主动落脚反馈必发散：人也靠调步宽稳横向]\label{ins:lb-lateral-must}
\cref{eq:lb-alip} 不只是一个工程拟合，它印证了一个被生物力学反复证实的事实：\textbf{横向平衡只能靠落脚，且主要靠调\emph{步宽}（而非步长）。}Kuo 与 Bauby 的人体行走研究\cite{kuo1999stabilization} 指出，人在矢状面（前后）相对稳定、横向（左右）才是更需要主动控制的方向，且人主要通过\emph{调节左右脚的横向落点（步宽）}来稳定横向——与矢状面靠调步长（捕获点 \cref{eq:lb-capture}）形成分工。把这条迁移到点足双足：若控制器\emph{缺少}主动的横向落脚反馈（例如只做了矢状捕获、没做横向的 \cref{eq:lb-alip}），roll 模态\emph{必然}发散——这正是工程实践中“站不住总是 roll 先倒”这一指纹现象的根（详见 \cref{ch:legged-debug} 对该指纹的复盘）。\textbf{矢状靠步长、横向靠步宽，是双足落脚平衡的二维分工。}
\end{insight}

\subsection{反应式 capture 层的 DCM 投影几何}
\label{subsec:lb-dcm-proj}

\cref{eq:lb-alip} 是“规划层”的横向落脚律。在“反应式”补偿层（即在 MPC 之外，根据实测状态实时算一个落脚目标去喂给 WBC 或参考），工程上常用一个更直接的 DCM 投影几何。它把当前 DCM 投影到落脚目标：

\begin{align}
  \mathbf{dcm}&=\mathbf{base}_{xy}+\frac{\mathbf{vel}_{world}}{\omega},\label{eq:lb-dcm-proj1}\\
  \mathbf{projected\_dcm}&=\mathbf{p}_{\mathrm{stance},xy}+\exp\!\big(\mathrm{clamp}(\omega\,t,\,0,\,1.5)\big)\cdot(\mathbf{dcm}-\mathbf{p}_{\mathrm{stance},xy}).\label{eq:lb-dcm-proj2}
\end{align}

\cref{eq:lb-dcm-proj1} 就是 \cref{eq:lb-dcm-def} 的向量、机体平面形式（用归一化线动量 \texttt{state[0:3]/\allowbreak mass} 作质心速度的近似亦可）。\cref{eq:lb-dcm-proj2} 是把 DCM 沿 \cref{eq:lb-dcm-sol} 的 $e^{\omega t}$ 律向前\emph{外推}一段时间 $t$，得到“若不干预，DCM 将走到的位置”，据此设落脚目标。

\begin{pitfall}[指数必须 clamp：$e^{\omega t}$ 在数值上会爆]\label{pit:lb-clamp-exp}
\cref{eq:lb-dcm-proj2} 里的指数项 $\exp(\omega t)$ 直接来自 DCM 的发散解 \cref{eq:lb-dcm-sol}——它\emph{在数学上}就是要指数增长的。但在代码里不加约束地算 $\exp(\omega t)$，当 $t$ 稍大（或求解器时序异常给出一个大的 $t$）时会迅速溢出到天文数字，把落脚目标投到几十米外，整条反应式 capture 层瞬间失控。故源码用 $\mathrm{clamp}(\omega t,0,1.5)$ 把指数自变量夹在 $[0,1.5]$（即外推时间不超过约 $1.5/\omega\approx 0.3\,\mathrm{s}$）。这是“理论上正确的发散律 + 工程上必须的数值护栏”的典型例子：发散是物理的真，但你不能让它在浮点里真的发散。
\end{pitfall}

\paragraph{文献定位：落脚律的谱系。}
横向/欠驱双足的落脚律是一个活跃的研究方向，本节给出的 ALIP 闭式律（\cref{eq:lb-alip}）只是其中一支。把它放进谱系（\cref{tab:lb-footstep-lit}）有助于理解工程取舍：所有这些方法都\emph{不}把捕获点当成无约束目标，而是结合可达性、步时、ZMP/接触可行性、角动量来用。

\begin{table}[H]\centering\small
\caption{欠驱/点足双足落脚律文献谱系}
\begin{tabular}{p{0.26\textwidth} p{0.30\textwidth} p{0.30\textwidth}}
\toprule
方法 & 核心思想 & 出处 \\
\midrule
ALIP（绕接触点角动量 LIP） & 用角动量描述 LIP，给横向交替步宽闭式律（本节 \cref{eq:lb-alip}） & Gong \& Grizzle\cite{gong2022alip} \\
ALIP-MPC & 把\emph{落脚}作为 MPC 的决策变量在 ALIP 上滚动求解 & Gibson/Gong/Grizzle\cite{gibson2022alipmpc} \\
H-LIP（Hybrid-LIP） & 混合系统 LIP，step-to-step 动力学 + 反馈综合稳定 & Xiong \& Ames\cite{xiong2022hlip} \\
DCM 步位 + 步时调整 & 同时调\emph{落脚位置}与\emph{步时}稳定 DCM & Khadiv 等\cite{khadiv2016stepping} \\
MPFC（落脚 MPC 谱系） & 把落脚约束嵌入 MPC 的统一框架 & Acosta \& Posa\cite{acosta2023mpfc} \\
捕获点框架（奠基） & 一步/N 步可捕获区域，push recovery & Pratt 等\cite{pratt2006capture} \\
\bottomrule
\end{tabular}
\label{tab:lb-footstep-lit}
\end{table}

\begin{practice}
\begin{enumerate}
  \item 用一句话解释为什么横向步宽 $W=0$ 会让 roll 发散，而矢状面没有这个“必须非零”的约束（提示：横向支撑点靠交替来接住发散，$W=0$ 则不再交替）。
  \item 在 \cref{eq:lb-alip} 中取 $\ell T\to 0$（步时极短）的极限，用 $\sinh(\ell T)\approx \ell T$、$\cosh(\ell T)\approx 1$ 估计 $L_x^{des}$ 的量级，说明步时趋零时所需横向落脚量如何变化。再取 $\ell T$ 较大时的行为做对比。
  \item （承接 \cref{sec:lb-lip}）矢状靠步长、横向靠步宽（\cref{ins:lb-lateral-must}）。把这二者合起来，写出“一步落脚目标”的二维结构：矢状分量用捕获点 \cref{eq:lb-capture}、横向分量用 \cref{eq:lb-alip}。这就是一个完整的反应式落脚规划器的骨架。
  \item （概念）\cref{eq:lb-dcm-proj2} 把 DCM 沿 $e^{\omega t}$ 外推。若把外推时间 $t$ 设为“到本步触地剩余时间”，这个投影的物理含义是什么？为什么这比直接用当前 DCM 作落脚目标更合理？
\end{enumerate}
\end{practice}

% ════════════════════════════════════════════════════════════════
\section{把硬约束接进 OCS2：五块拼图与 ReferenceManager}
\label{sec:lb-ocs2}

\paragraph{动机：从“最小模型”回到真实求解栈。}
\cref{sec:lb-lip,sec:lb-alip} 是“为什么迈步能稳”的闭式底座。但真正跑在机器人上的不是这些闭式律，而是 \cref{ch:quad_mpc} 那套 OCS2 质心 NMPC——它把落脚、力分配、步态时序都交给一个在线滚动的非线性最优控制问题（OCP）去解。把四足栈移植到点足，要在这个 OCP 里改两类东西：把\emph{对点足才对}的参数（摩擦系数、权重、站高）配对，以及给它\emph{接入新的约束}（例如更严的法向力下界）。本节先建立 OCS2 的架构心智模型（不建立它，任何在线改 OCP 的尝试都会“改了没反应”或“与求解器打架”），再给出硬约束接入的清单与三个静默陷阱。

\subsection{OCP 五块拼图与静态/动态分界线}
\label{subsec:lb-five-pieces}

要把一个抽象的 OCP（\cref{ch:quad_mpc} 的 \cref{sec:qp,sec:state-space} 已建立其离散 multiple-shooting 形式）落到一台具体机器人上，OCS2 工程上把它拆成\textbf{必须分别提供的五块拼图}（\cref{tab:lb-five-pieces}）：动力学、代价、约束、参考、模式调度。装配入口（源码里的 \texttt{setupOptimalControlProblem}）的全部工作，就是把 \texttt{task.info} / URDF / \texttt{reference.\allowbreak info} 三个配置文件“编译”成一个 \texttt{OptimalControlProblem} 容器 + 一个 \texttt{ReferenceManager}——它本身\emph{不求解、不下发}。

\begin{table}[H]\centering\small
\caption{OCS2 的 OCP 五块拼图与静态/动态分界（\cref{ch:quad_mpc} 的 OCP 在工程上的落地形态）}
\begin{tabular}{clll}
\toprule
\# & 拼图 & 内容 & 静/动态 \\
\midrule
\cnum{1} & 动力学 $f$ & 质心模型 $x_{k+1}=f_d(x_k,u_k)$ & 静态：挂 \texttt{problemPtr\_} \\
\cnum{2} & 代价 $\ell/\Phi$ & 多项叠加到 \texttt{costPtr->add} & 静态：挂 \texttt{problemPtr\_} \\
\cnum{3} & 约束 $g/h$ & 逐脚 add，分等式/不等式/软约束 & 静态：挂 \texttt{problemPtr\_} \\
\cnum{4} & 参考 $x_{ref}$ & \texttt{/cmd\_vel}$\to$TargetTrajectories & \textbf{动态：进 \texttt{referenceManagerPtr\_}} \\
\cnum{5} & 模式调度 ModeSchedule & gait（谁何时触地） & \textbf{动态：进 \texttt{referenceManagerPtr\_}} \\
\bottomrule
\end{tabular}
\label{tab:lb-five-pieces}
\end{table}

\textbf{分界线就是 ReferenceManager}：前三块（动力学/代价/约束）相对\emph{静态}——开机装配一次、运行中基本不变，直接挂在 \texttt{problemPtr\_} 上；后两块（参考/模式调度）是\emph{动态}的——随指令在线变，于是\emph{不挂} problem，而封装进 \texttt{referenceManagerPtr\_}。

\begin{insight}[设计因果链：为了并行 $\Rightarrow$ OCP 必须可拷贝 $\Rightarrow$ 不能直接改 problem $\Rightarrow$ 需 ReferenceManager]\label{ins:lb-refmgr-why}
为什么非要分这条线、非要 ReferenceManager 这个看似多余的间接层？这是移植/扩展 OCS2 前\emph{必须}建立的心智模型，它的因果链是：
\begin{enumerate}
  \item \textbf{为了并行}：SQP/DDP 求解器用多线程做 rollout / 线性化，\emph{每个工作线程持有一份 OCP 的拷贝}，以换取无锁的安全并行；
  \item \textbf{$\Rightarrow$ OCP 必须可自由拷贝}：于是 OCP 被设计成“可被自由拷贝的纯问题描述”；
  \item \textbf{$\Rightarrow$ 不能直接改 problem}：后果有二——(a) 你在外面改原始 \texttt{problemPtr\_} 的参考值，求解器线程用的是\emph{自己的副本}、根本看不到，\emph{改了等于没改}；(b) 即便共享字段，前台写、后台多线程读 $=$ 数据竞争，且可能拼出“半新半旧”的不自洽问题致策略发散；
  \item \textbf{$\Rightarrow$ 需要另一条受框架管理、线程安全、所有副本都能看到的更新通道}：这就是 \texttt{ReferenceManager + SolverSynchronizedModule}。其落地是\textbf{三段式}——ROS 回调线程只加锁\emph{缓冲} $\to$ 求解前回调 \texttt{preSolverRun} 安全\emph{写入} $\to$ 求解器 \texttt{modifyReferences} \emph{查询}。
\end{enumerate}
一句话：\textbf{任何“在线改 OCS2 行为”的功能（换步态、改参考、注入落脚 XY）都必须套这个三段式经 ReferenceManager 灌入}；在 post-MPC 的 WBC 层注入、或把动态参考焊死进静态代价里，都违背这条分界 $\Rightarrow$ 改了没反应、或与 MPC 参考打架。这条洞见也直接解释了 \cref{sec:lb-runtime} 里“运行时注入 gait 有 $\sim 1$ 个时域延迟”的现象。
\end{insight}

\subsection{点足才对的真实配置：别照抄四足值}
\label{subsec:lb-config}

移植中一个反复出事的点是：大量文档/默认配置里的数值是\emph{四足}的，对点足是\emph{错的}。\cref{tab:lb-braver-config} 列出点足（Braver）相对四足的关键配置差异，每一项都有“点足为什么必须这样”的理由。

\begin{table}[H]\centering\small
\caption{点足双足（Braver）相对四足的关键 OCS2 配置（对点足才对的值）}
\begin{tabular}{p{0.30\textwidth} p{0.24\textwidth} p{0.38\textwidth}}
\toprule
参数 & 点足值（vs 四足） & 为什么 \\
\midrule
\texttt{frictionCoefficient} & $0.6$（四足常 $0.3$） & 须 $\le$ 仿真地面 $\mu$；且 OCS2 软锥与 WBC 金字塔两层\emph{必须一致} \\
\texttt{comHeight} & $0.42737\,\mathrm{m}$ & 即 LIP 高度 $h$，定 $\omega=\sqrt{g/h}\approx 4.8$ \\
\texttt{sqpIteration} & $3$（四足常 $1$） & 点足更难，靠 warm-start 多迭代 \\
\texttt{timeHorizon} & $1.0\,\mathrm{s}$ & 须 $\ge$ 一个步态周期（$\approx 0.6\,\mathrm{s}$） \\
\texttt{Q[6:12]}（base 位姿段） & $[1000,1000,2000,100,400,400]$ & \textbf{z=2000 全表最大}=守站高；pitch/roll=400 刻意压姿态（点足无静稳） \\
\texttt{defaultModeSequenceTemplate} & \texttt{standing\_trot}（非 \texttt{DOUBLE\_SUPPORT}） & 必须默认就踏步，见 \cref{sec:lb-runtime} \\
\bottomrule
\end{tabular}
\label{tab:lb-braver-config}
\end{table}

\paragraph{为什么 base 位姿权重 z=2000 是全表最大？}
回到 \cref{ins:lb-no-static}：点足没有静稳，机身高度 $z$ 与姿态 pitch/roll 全靠 MPC“摁住”。\cref{tab:lb-braver-config} 里 \texttt{Q[6:12]} 的 $z$ 分量取 $2000$（全状态最大），就是在代价里给“维持站高”最高优先级；pitch/roll 取 $400$（也偏大）压姿态。这与四足“给腿留自由度、让 MPC 自己找 gait”的松权重哲学形成对比——四足有多边形兜底，可以把状态权重放松；点足必须把状态摁得很紧。这是“点足无静稳”这一物理事实在\emph{权重}层面的直接体现。

\subsection{硬不等式约束的接入清单}
\label{subsec:lb-hard-constraint}

把一条新的\emph{硬不等式}约束（例如要求 stance 脚法向力 $f^z\ge 50\,\mathrm{N}$，强制不抬脚）接进 OCS2 的 OCP，标准步骤如下：

\begin{codebox}[C++]{OCS2 接入硬不等式约束的最小骨架（\cref{ch:quad_mpc} OCP 的 inequalityConstraint 容器）}
// 1. 继承 StateInputConstraint，声明约束类型为不等式
class FootNormalForceConstraint : public ocs2::StateInputConstraint {
public:
  FootNormalForceConstraint(/* legIdx, floor */)
    : ocs2::StateInputConstraint(ocs2::ConstraintOrder::Linear) {}  // 声明 ConstraintOrder

  // 2. 何时激活:只在该脚 stance 时（务必重写，默认 true 会在错误模式生效）
  bool isActive(ocs2::scalar_t t) const override { return isStanceLeg(t, legIdx_); }

  // 3. 约束维数（首次构建曾因漏写此函数编译失败）
  size_t getNumConstraints(ocs2::scalar_t /*t*/) const override { return 1; }

  // 4. 约束值 h(x,u) >= 0:这里 h = f_z - floor
  ocs2::vector_t getValue(ocs2::scalar_t t, const ocs2::vector_t& x,
                          const ocs2::vector_t& u, const ocs2::PreComputation&) const override {
    ocs2::vector_t h(1);
    h(0) = contactForceZ(u, legIdx_) - floor_;   // f_z - 50
    return h;
  }

  // 5. 线性近似（不等式约束的 Jacobian）
  ocs2::VectorFunctionLinearApproximation getLinearApproximation(/* ... */) const override { /* ... */ }
};
// 6. 在 OCP 里按脚注册到 inequalityConstraintPtr
problem.inequalityConstraintPtr->add("footNormalForce_LF", std::move(constraintLF));
\end{codebox}

要点：继承 \texttt{StateInputConstraint} $\to$ 在构造里声明 \texttt{ConstraintOrder} $\to$ 实现 \texttt{isActive()} / \texttt{getNumConstraints()} / \texttt{getValue()} / \texttt{getLinearApproximation()} 四个虚函数 $\to$ 在 OCP 的 \texttt{inequalityConstraintPtr} 上\emph{按脚}注册。

\subsection{三个静默陷阱：注册不等于满足，命名不等于含义}
\label{subsec:lb-traps}

接入硬约束\emph{看起来}做完了，但有三个静默陷阱会让你误以为约束生效，实则不然。

\begin{pitfall}[陷阱一：注册硬约束 $\ne$ 闭环满足，\texttt{g\_max} 收敛不含不等式残差]\label{pit:lb-register-ne-satisfy}
你注册了硬约束 $f^z\ge 50\,\mathrm{N}$，求解器报告“收敛”（\texttt{sqp.g\_max} 达标），你以为大功告成。\textbf{错。}OCS2 本地的 \texttt{totalConstraintViolation} 只算 $\sqrt{\texttt{dynamicsSSE}+\texttt{equalitySSE}}$，\emph{不含}不等式残差 SSE。也就是说，\texttt{g\_max} 收敛\emph{不}保证硬不等式被 active policy 满足。判断硬不等式是否真被满足的\textbf{唯一}求解器级证据，是单独去读 \texttt{getPerformanceIndices().\allowbreak inequalityConstraintsSSE}。实测把这点砸实：给 hard \texttt{Fz>=50N} 时，\texttt{inequalityConstraintsSSE} 高达 $66680$、\texttt{dynamicsViolationSSE} 反而从 $3.59$ 涨到 $66.38$（硬约束\emph{反而引入}更大动态不一致）、当前 floor violation $1633\,\mathrm{N}$、policy horizon violation $1358\,\mathrm{N}$（遍布整条 horizon 而非仅当前采样点）。结论：\textbf{力不可行是机身失稳的\emph{后果}（身体倾倒 $\Rightarrow$ 为追踪漂移的参考而要求不可执行的接触力），不是摩擦不够；在不可行参考上加硬约束、加迭代，只会让 desired force 爆得更大。}（这条“指标改善 $\ne$ 闭环改善”的诊断学，是 \cref{ch:legged-debug} 的核心方法论之一。）
\end{pitfall}

\begin{pitfall}[陷阱二：\texttt{frictionConeSoftConstraint.\allowbreak mu} 是惩罚权重，不是摩擦系数]\label{pit:lb-mu-naming}
软摩擦锥约束的配置里有一个字段 \texttt{frictionConeSoftConstraint.\allowbreak mu}（常默认 $1.0$）。新人极易把它当成\emph{摩擦系数}去调大，以为这样能让机器人“更不容易打滑”。\textbf{它根本不是摩擦系数}——它是 \texttt{RelaxedBarrierPenalty}（松弛障碍惩罚）的\emph{惩罚权重}，配合 \texttt{delta}（软化半径）控制“越红线罚多重”。真正的摩擦系数是另一个字段 \texttt{frictionCoefficient}。把 \texttt{mu} 调大不会改变摩擦可行集（红线位置没变），反而把软约束惩罚调重、使 Hessian 更刚、求解病态。记忆口诀：\textbf{改“红线在哪”动 \texttt{frictionCoefficient}；改“越红线罚多重”才动 \texttt{mu}/\allowbreak\texttt{delta}}。
\end{pitfall}

\begin{pitfall}[陷阱三：默认软圆锥不含 $f^z\ge 0$，负法向力“免费”]\label{pit:lb-soft-cone-free}
OCS2 的软摩擦锥默认形如 $\mu(f^z+f_{\mathrm{grip}})-\sqrt{(f^x)^2+(f^y)^2+\texttt{reg}}\ge 0$，且默认正则 \texttt{regularization}$=25$（软锥要到 $f^z\approx 8.3\,\mathrm{N}$ 才真正到边界）。关键缺陷：\textbf{这个软锥不含独立的 $f^z\ge 0$ 项}——也就是说，对软锥而言，\emph{负}的法向力是“免费”的（不受惩罚）。这就解释了一个困惑现象：OCS2 解出的 desired stance 力会频繁出现\emph{负} $f^z$（脚“拉”地面）。这不只是身体倒导致的，而是软锥本身对负 $f^z$ 无惩罚。理解这点，你才不会误把“negative stance Fz”全归因于身体失稳，也才知道若要禁止负法向力，得\emph{另外}显式加一条 $f^z\ge 0$（或更高的 floor）的约束——但要小心陷阱一（注册 $\ne$ 满足）。
\end{pitfall}

\begin{practice}
\begin{enumerate}
  \item 用 \cref{ins:lb-refmgr-why} 的因果链解释：若有人为了“运行时改落脚目标”，直接在 ROS 回调里修改 \texttt{problemPtr\_} 里的某个参考字段，会发生什么？（提示：求解器线程用的是哪份拷贝？）
  \item 你给 OCS2 注册了硬约束 $f^z\ge 50\,\mathrm{N}$，跑起来求解器报“收敛”。写出你接下来\emph{必须}检查的那个标量字段名，并说明为什么 \texttt{g\_max} 收敛不足以下结论（\cref{pit:lb-register-ne-satisfy}）。
  \item 同事说“我把 \texttt{frictionConeSoftConstraint.\allowbreak mu} 从 $1.0$ 调到 $3.0$，想让机器人更不容易打滑”。指出这句话的两个错误，并告诉他正确该动哪个字段（\cref{pit:lb-mu-naming}）。
  \item （概念）\cref{tab:lb-braver-config} 给 base 位姿 z 权重取全表最大 $2000$。把它和 \cref{ins:lb-no-static}“点足无静稳”联系起来：为什么四足可以把状态权重放松、让腿留自由度，而点足必须把 z/pitch/roll 摁紧？
\end{enumerate}
\end{practice}

% ════════════════════════════════════════════════════════════════
\section{点足 WBC：把机身任务从“纯前馈”拆分出来}
\label{sec:lb-wbc}

\paragraph{动机：本章的“题眼”——一个被悄悄置零的反馈增益。}
到这里，矢状/横向的落脚理论（\cref{sec:lb-lip,sec:lb-alip}）和 OCS2 的硬约束/配置（\cref{sec:lb-ocs2}）都对了，但机器人\emph{还是}站不住、总是 roll 先倒、力矩撞饱和。\cref{sec:lb-ocs2} 的陷阱一已埋下伏笔：力不可行是\emph{后果}不是病根。真正的病根在\emph{执行层}——WBC 的机身（base）加速度任务被替换成了四足的\emph{纯前馈}版本，机身位姿/高度/角度的反馈增益\emph{恒为 $0$}。四足靠多边形兜底能活（\cref{pit:lb-quad-intuition}），点足是倒立摆、掐断这条唯一的姿态反馈通道必然发散。本节给出正确的拆分形式，并解析旧前馈写法\emph{错在哪}。

\subsection{机身加速度任务的拆分形式}
\label{subsec:lb-base-split}

回忆 \cref{ch:quad_wbc} 的浮动基 WBC 子任务体系（\cref{sec:qw-floating}，\cref{tab:qw-subtasks}）：其中有一项是\emph{机身加速度任务}，它从质心动量算出期望的机身加速度，再写成对广义加速度的等式任务。四足的实现是\emph{纯前馈}的（只用质心动量算期望加速度，不加任何位姿/速度反馈）。点足的正确实现必须把它\textbf{拆分}成三块带反馈的子任务：

\begin{equation}
  \underbrace{\texttt{formulateBaseAccelTask}}_{\text{机身加速度任务}}
  =\underbrace{\texttt{BaseXYLinearAccel}}_{\text{XY 线加速度(前馈)}}
  +\underbrace{\texttt{BaseHeightMotion}}_{\text{高度(Jacobian+PD)}}
  +\underbrace{\texttt{BaseAngularMotion}}_{\text{角(Jacobian+旋转误差+世界系}\boldsymbol\omega)}.
  \label{eq:lb-base-split}
\end{equation}

其中\textbf{角任务}是最关键、也最容易写错的一块。它的正确形式是：任务雅可比取机身角部分 $\mathbf{a}=\texttt{base\_j}_{(3:6)}$，期望右端

\begin{equation}
  \mathbf{b}=\mathbf{a}_{des}
  +\mathbf{K}_p\cdot\underbrace{\texttt{rotationErrorInWorld}(\mathbf{R}_{des},\mathbf{R}_{meas})}_{\text{世界系旋转矩阵误差}}
  +\mathbf{K}_d\cdot(\boldsymbol\omega_{des}-\boldsymbol\omega_{meas})
  -\texttt{base\_dj}_{(3:6)}\cdot\mathbf{v}.
  \label{eq:lb-base-angular}
\end{equation}

\cref{eq:lb-base-angular} 里的期望量 $\mathbf{a}_{des},\boldsymbol\omega_{des}$（以及高度块的期望）由 \texttt{computeBaseKinematicsFromCentroidalModel} 从质心模型\emph{解析}地算出，\textbf{不做有限差分}（理由见 \cref{pit:lb-finite-diff}）。修复后的 \texttt{task.info} 把这三块的增益从 $0$ 打开：\texttt{baseHeight/\allowbreak baseAngular} 的 $K_p:0\to20$、$K_d:0\to3$，关节 $K_d:0\to3$。

\begin{insight}[拆分 base 任务 = 浮动基子任务在“点足无静稳”下的特化]\label{ins:lb-split-is-specialization}
\cref{eq:lb-base-split} 不是新发明，而是 \cref{sec:qw-floating} 那套浮动基子任务（\cref{tab:qw-subtasks}）在点足上的\emph{特化}。四足那一项机身加速度任务\emph{也}本可以带反馈，只是四足用纯前馈\emph{就够了}——因为支撑多边形会兜底（\cref{ins:lb-no-static}）。点足把这层兜底撤掉后，原本“可省”的反馈就变成“续命”的（\cref{pit:lb-quad-intuition}）。所以本节的工作可以一句话概括：\textbf{把 \cref{tab:qw-subtasks} 里那个对四足够用的纯前馈机身任务，在点足上拆成“XY 线加速度（仍前馈）+ 高度（PD）+ 角（旋转误差 PD）”三块带反馈的子任务}，其中角任务必须用世界系旋转误差（下一小节解释为什么不能用欧拉角直接相加）。这也是“fork WBC 时必审 base 任务反馈增益”这条移植铁律的来源——四足静稳会\emph{掩盖}增益为 $0$ 的缺陷，点足倒立摆不会。
\end{insight}

\subsection{旧前馈写法的两块根因}
\label{subsec:lb-two-bugs}

为什么旧的纯前馈写法不只是“缺反馈”、而是\emph{藏着错误}？它有两块互相独立的根因拼图。

\begin{pitfall}[根因一（frame bug）：把欧拉角误差直接加到机体系角加速度 = 三切空间相加]\label{pit:lb-frame-bug}
旧的角反馈（若有人试着加一点）常这样写：把 ZYX 欧拉角的误差、加上欧拉角速率误差，\emph{直接加到机体系的角加速度} $\mathbf{b}.\texttt{tail}\langle3\rangle()$ 上。\textbf{这是错的}，因为它把\emph{三个不同切空间}的量当成同类相加：欧拉角误差活在欧拉角的局部坐标里、机体系角加速度活在机体系切空间、世界系量又是另一个——它们之间差着依赖姿态的映射（\cref{ch:lie}），不能直接相加。后果是近 gimbal（万向锁）奇异、跨轴耦合：小角度（近直立）时\emph{可能}凑合，大角度必错。这正是工程实践中“单独打开角反馈反而让站立更糟”的精确机理——它让人误以为“反馈有害”，从而强化“保持增益 $=0$”的错误结论。正解必须用 \cref{eq:lb-base-angular} 的 \texttt{rotationErrorInWorld}（世界系旋转矩阵误差，本质是 $\log(\mathbf{R}_{des}\mathbf{R}_{meas}^\top)$ 的轴角，\cref{ch:lie}）+ 世界系 $\boldsymbol\omega$ + base 角 Jacobian，把误差、角速度、加速度都统一到\emph{同一个}切空间再相加（与 \cref{sec:qw-practice} 的欧拉角误差变基同源）。
\end{pitfall}

\begin{pitfall}[根因二（有限差分噪声放大 $\sim 476\times$）：为何切到 FullCentroidal 才致命]\label{pit:lb-finite-diff}
旧前馈写法里还有一处隐患：它用有限差分算关节加速度 $\texttt{jointAccel}=(u-u_{\mathrm{last}})/\texttt{period}$，其中 $\texttt{period}\approx 0.0021\,\mathrm{s}$（实时控制环周期）做\emph{分母}。把逐拍的输入抖动除以这么小的数，等于把噪声放大约 $1/0.0021\approx 476$ 倍。为什么这个雷在四足/SRBD 下\emph{没爆}？因为在 SRBD（单刚体）质心模型下，机身动力学对关节加速度的耦合项 $\mathbf{A}_j\approx 0$（SRBD 用名义惯量预算一次、忽略关节动量贡献，\cref{ch:quad_mpc} 的 \texttt{centroidalModelType}），这个被放大 $476\times$ 的噪声项乘上 $\approx 0$ 而\emph{幸免}。\textbf{但一旦切到 FullCentroidal（完整质心模型，$\mathbf{A}_j\ne 0$）就致命}——放大的噪声直接灌进机身任务。这解释了 \texttt{centroidalModelType} 选择与机身任务噪声的耦合，也是 \cref{eq:lb-base-angular} 强调“期望量须\emph{解析}（\texttt{computeBaseKinematicsFromCentroidalModel}）、不做有限差分”的原因。
\end{pitfall}

\begin{remark}[修复后是欠阻尼过冲，不是符号错]\label{rem:lb-underdamped}
按 \cref{eq:lb-base-split,eq:lb-base-angular} 修复后，单纯站立\emph{仍}会在约 $1.3\,\mathrm{s}$ 后倒——但这次是\emph{良性}的：onset trace 显示反馈\emph{已起效}（pitch 初始后倒被拉回、过零，再前向发散），这是典型的\textbf{欠阻尼过冲}而非符号错。佐证：此时\emph{加大} base 权重反而最快倒（与欠阻尼一致，权重越高越容易过冲）。于是问题从“反馈方向错/缺反馈”这种结构性 bug，\emph{降级}成了一个阻尼调参问题。结合“点足无静稳、站立必须是踏步”（\cref{ins:lb-no-static}），最终的解不是把站立调到无限久，而是让它\emph{从一开始就踏步}（\cref{sec:lb-runtime}）——一旦机身反馈正确、再叠加确定性周期 gait + 纯 MPC 落脚，机器人就稳定走起来了（这也印证了 \cref{rem:lb-implicit-capture}：MPC 隐式在做捕获点的事）。
\end{remark}

\begin{practice}
\begin{enumerate}
  \item 用 \cref{ins:lb-split-is-specialization} 的语言说明：四足的纯前馈机身任务和点足的拆分机身任务，\emph{多刚体动力学}上是同一项，差别只在“是否带反馈”。为什么这个差别在四足上无所谓、在点足上是生死线？
  \item \cref{pit:lb-frame-bug} 说欧拉角误差不能直接加到机体系角加速度。用 \cref{ch:lie} 的 $\dot{\mathbf{R}}=\boldsymbol\omega^\wedge\mathbf{R}$ 与“旋转误差应取 $\log(\mathbf{R}_{des}\mathbf{R}_{meas}^\top)$”说明：为什么 \texttt{rotationErrorInWorld} 才是把误差放进\emph{正确切空间}的做法？为什么近直立小角度时旧写法“可能凑合”？
  \item 估算 \cref{pit:lb-finite-diff} 的噪声放大倍数：$\texttt{period}=0.0021\,\mathrm{s}$，有限差分 $(u-u_{\mathrm{last}})/\texttt{period}$ 的放大系数是多少？再解释为什么 SRBD（$\mathbf{A}_j\approx 0$）幸免、FullCentroidal（$\mathbf{A}_j\ne 0$）致命。
  \item （承接 \cref{rem:lb-underdamped}）修复后站立仍 $\sim 1.3\,\mathrm{s}$ 倒，但是\emph{欠阻尼过冲}。给出两条证据区分“欠阻尼”与“符号错/缺反馈”（提示：onset trace 的 pitch 是否先被拉回过零？加大 base 权重会更稳还是更快倒？）。
\end{enumerate}
\end{practice}

% ════════════════════════════════════════════════════════════════
\section{运行时步态注入与“从第一拍就走”}
\label{sec:lb-runtime}

\paragraph{动机：机身反馈对了、落脚理论对了，最后一步是“怎么让它一开机就踏步”。}
\cref{rem:lb-underdamped} 已点明：点足的“站立”在物理上只能是“踏步”（\cref{ins:lb-no-static}）。于是控制器\emph{不能}先进入一个静态双支撑站立态、再等指令切换到踏步——它必须\emph{从第一个 MPC policy 就规划踏步}。这看似简单，却卡在 OCS2 步态注入的一个时序细节上，并连带一个新手最常崩的步态序列陷阱。

\subsection{运行时注入 gait 的约一个时域延迟}
\label{subsec:lb-gait-delay}

直觉做法是：开机先双支撑站着，等需要走时通过运行时接口（源码里的 \texttt{GaitReceiver}）注入一个踏步序列。但 \texttt{GaitReceiver} 是在当前时域的\emph{末端}（\texttt{finalTime}）才把新序列 \texttt{insert} 进去的——于是注入的 gait 有\textbf{约 $1$ 个时域（$\sim 1\,\mathrm{s}$）的延迟}才生效。对一个没有静稳的倒立摆，这 $1\,\mathrm{s}$ 的“干等双支撑”足以让它倒掉。这个延迟正是 \cref{ins:lb-refmgr-why} 那条设计因果链的直接推论——运行时改 gait 走的是 ReferenceManager 的同步边界通道，而它在时域末端才安全写入。

\textbf{正解}：不靠运行时注入，而是直接改 \texttt{reference.\allowbreak info} 里的 \texttt{defaultModeSequenceTemplate}——把它从 \texttt{DOUBLE\_SUPPORT}（双支撑站立）改成 \texttt{standing\_trot}（带踏步的序列）。这样\emph{第一个} MPC policy 就规划的是踏步，机器人一开机就走，没有那 $1\,\mathrm{s}$ 的致命空窗（\cref{tab:lb-braver-config} 已列此项）。

\subsection{ModeSequenceTemplate 平铺陷阱}
\label{subsec:lb-tile-trap}

定义步态序列时有一个新手最常崩的坑，源于 OCS2 把\emph{模板}平铺（tile）成完整时域序列的方式（\texttt{tileModeSequenceTemplate}，对照 \cref{ch:quad_gait} 的步态调度器 \cref{sec:qg-scheduler}）。

\begin{pitfall}[平铺陷阱：单模板被铺成连续单支撑，且切换时刻数必比模式数多一个]\label{pit:lb-tile}
若你天真地把模板写成单个模式 \texttt{[RIGHT\_STANCE]}（想表达“右脚支撑”），\texttt{tileModeSequenceTemplate} 会把它平铺成\emph{连续的}右支撑——结果左脚\emph{整个时域都在 swing}（一直抬着不落），完全不是交替步态。\textbf{双足的一次性步必须把模板写成 \texttt{stance$\to$DS$\to$opposite\_stance$\to$DS}}（单支撑$\to$双支撑过渡$\to$对侧单支撑$\to$双支撑），才能平铺成正确的左右交替步态。

更隐蔽的崩溃源是\textbf{长度对不上}：在 OCS2 的步态模板里，\texttt{switchingTimes}（切换时刻数组）\emph{必须比} \texttt{modeSequence}（模式数组）\textbf{多一个元素}。这是因为 $N$ 个模式之间有 $N+1$ 个时刻边界（首时刻 + 每个模式的结束时刻）。新手定义步态时最常见的崩溃就是“切换时刻数和模式数一样多、少了一个”，导致平铺逻辑越界或最后一个模式“永久持续”。代码级根因：当模板为空时 \texttt{tileModeSequenceTemplate} 会直接 \texttt{return}（致“最后一个模式永久持续”），平铺循环 \texttt{while eventTimes.\allowbreak back()<finalTime} 整体重复贴模板，而每段时长 \texttt{deltaTime=templateTimes[i+1]-templateTimes[i]} 正是从“时刻数比模式数多一个”这个结构来的。
\end{pitfall}

\begin{insight}[同步边界只解决“谁何时改”，不保证落脚动力学合理]\label{ins:lb-sync-boundary}
关于步态/落脚的运行时更新，有一条容易混淆的边界：OCS2 的\emph{同步边界}（经 ReferenceManager / SolverSynchronizedModule，\cref{ins:lb-refmgr-why}）是 target/mode 更新的\emph{正确位置}——它解决的是“ownership（谁拥有这个更新）”与“timing（何时安全写入）”，避免前后台数据竞争、避免与 MPC 参考打架。但它\emph{只}解决这两件事，\textbf{不保证你注入的落脚目标在动力学上是合理的}。也就是说，把落脚目标经正确通道、在正确时刻灌进去，不等于这个落脚点能让机器人稳——后者要靠 \cref{sec:lb-lip,sec:lb-alip} 的落脚律或 MPC 自己的优化。把“通道正确”与“内容正确”分开，是避免“我明明经 ReferenceManager 灌了落脚目标，怎么还是倒”这类困惑的关键。这也呼应了本章的最终方案：与其手写落脚启发式经同步边界灌进去，不如修好机身反馈（\cref{sec:lb-wbc}）后\emph{放手让 MPC 自己拥有落脚 XY}（\cref{rem:lb-implicit-capture}）。
\end{insight}

\paragraph{把指令变成 MPC 参考：\texttt{cmd\_vel} 接口。}
最后补一块拼图：用户的速度指令 \texttt{/cmd\_vel} 怎么变成 \cref{tab:lb-five-pieces} 第 \cnum{4} 块的参考轨迹？源码的 \texttt{cmdVelToTargetTrajectories} 取当前机身位姿 \texttt{state[6:12]}，把机体系速度指令旋到世界系 $\mathbf{v}_{world}=\mathbf{R}_{zyx}\,\mathbf{cmd}[0:3]$，再外推目标位置 $\mathbf{p}_{target}=\mathbf{p}_{cur}+\mathbf{v}_{world}\cdot t_{toTarget}$、目标偏航 $\psi_{target}=\psi_{cur}+\omega_z\cdot t_{toTarget}$，并把 $z$ 设为 \texttt{comHeight}、pitch/roll 设为 $0$、归一化动量段 $[0:6]$ 填 $\mathbf{v}_{world}$。这条参考再经 ReferenceManager 的三段式（\cref{ins:lb-refmgr-why}）安全灌入求解器。最终效果：机身反馈修复 + 默认踏步 gait + 纯 MPC 落脚（\texttt{swingFootTargetEnabled=false}，无外部 XY 目标、capture 层全关），机器人 $18\,\mathrm{s}$ 不倒、速度跟踪误差 $1\mathrm{e}{-4}$ 量级、走直线（$3.85\,\mathrm{m}$ 内横漂 $<2\,\mathrm{cm}$）。\textbf{这印证了本章的主线判断：那一整套复杂的 capture 踏步补偿，不过是给坏掉的机身任务打的补丁；根因（\cref{sec:lb-wbc}）一修，最简单的路线就通了。}

\begin{practice}
\begin{enumerate}
  \item 用 \cref{ins:lb-refmgr-why} 解释为什么运行时注入 gait 会有 $\sim 1$ 个时域延迟，以及为什么改 \texttt{defaultModeSequenceTemplate} 能绕开它（提示：注入走同步边界、在时域末端写入；默认模板在装配时就生效）。
  \item 写出双足“一次性步”的正确模式模板序列（\cref{pit:lb-tile}），并说明若误写成单个 \texttt{[RIGHT\_STANCE]} 会发生什么。
  \item 给定一个有 $4$ 个模式的步态序列，\texttt{switchingTimes} 该有几个元素？解释这个“$N+1$”从哪来（\cref{pit:lb-tile}）。
  \item （概念）\cref{ins:lb-sync-boundary} 区分“通道正确”与“内容正确”。举一个例子：你经 ReferenceManager 在正确时刻灌入了一个落脚目标，但机器人仍倒。这是“通道”的问题还是“内容”的问题？该去查 \cref{sec:lb-lip,sec:lb-alip} 的哪一部分？
\end{enumerate}
\end{practice}

% ════════════════════════════════════════════════════════════════
\section*{本章小结}
\addcontentsline{toc}{section}{本章小结}

本章把“四足栈移植到欠驱点足双足”的核心差异、公式底座与两个工程深坑串成一条线。一句话主线：\textbf{四足靠支撑多边形“接住”了很多本该由反馈完成的工作；点足把这层缓冲撤掉（支撑多边形退化成一个点、无静稳），于是平衡只能靠迈步，而那些被默认关掉、被纯前馈化的环节会逐一暴露。}

\begin{itemize}
  \item \textbf{点足无静稳}（\cref{sec:lb-pointfoot}）：两脚同 $x$ 线 $\Rightarrow$ 俯仰零力臂 $\Rightarrow$ 结构性地没有静止站立稳态，它是主动不稳的倒立摆，“站立”=“踏步”。把四足“base 反馈可选”的直觉搬过来是致命陷阱。
  \item \textbf{LIP/DCM/捕获点}（\cref{sec:lb-lip}）：$\ddot x=\omega^2 x$ 的\emph{正号}给出发散模态；DCM $\xi=x+\dot x/\omega$ 把不稳定性浓缩成一维 $\dot\xi=\omega(\xi-p)$；捕获点 $p=\xi$ 一步停稳。四足 gait 章的 $\sqrt{z_0/g}(\mathbf v-\mathbf v_{des})$ 就是捕获点反馈项。
  \item \textbf{ALIP 横向闭式律}（\cref{sec:lb-alip}）：横向平均速度为零却须以非零步宽 period-2 极限环行走；交替步宽 $L_x^{des}=\pm\frac12 mHW\,\frac{\ell\sinh(\ell T)}{1+\cosh(\ell T)}$；矢状靠步长、横向靠步宽。
  \item \textbf{OCS2 硬约束接入}（\cref{sec:lb-ocs2}）：五块拼图 + ReferenceManager 的因果链（为了并行$\Rightarrow$可拷贝$\Rightarrow$不能直接改 problem$\Rightarrow$需同步通道）；注册$\ne$满足（唯一证据 \texttt{inequalityConstraintsSSE}）；\texttt{mu} 是惩罚权重非摩擦系数；软锥对负 $f^z$ 免费。
  \item \textbf{点足 WBC 拆分}（\cref{sec:lb-wbc}）：机身任务从纯前馈拆成 XY 线加速度 + 高度 PD + 角（世界系旋转误差 PD）；旧前馈两块根因——切空间错配的 frame bug、有限差分噪声放大 $\sim 476\times$（SRBD 幸免、FullCentroidal 致命）。
  \item \textbf{运行时步态}（\cref{sec:lb-runtime}）：改 \texttt{defaultModeSequenceTemplate} 才能从第一拍就走；平铺陷阱（单模板铺成连续单支撑、切换时刻数比模式数多一个）。
\end{itemize}

\begin{table}[H]\centering\small
\caption{本章主要符号}
\begin{tabular}{lll}
\toprule
符号 & 含义 & 出处 \\
\midrule
$\omega=\sqrt{g/h}$ & LIP 自然频率（矢状），Braver $\approx 4.8$ & \cref{eq:lb-lip} \\
$\ell$ & 横向 LIP 频率 & \cref{eq:lb-alip} \\
$h$ ($z_0$) & 质心高度（恒定，LIP 假设），$\approx 0.427\,\mathrm{m}$ & \cref{eq:lb-lip} \\
$\xi=x+\dot x/\omega$ & 发散分量 DCM / 瞬时捕获点 & \cref{eq:lb-dcm-def} \\
$p_{\mathrm{capture}}=\xi$ & 捕获点（一步停稳的落脚点） & \cref{eq:lb-capture} \\
$W$ & 步宽（横向左右脚间距） & \cref{eq:lb-alip} \\
$T$ & 步时（一步时长） & \cref{eq:lb-alip} \\
$L_x^{des}$ & 横向交替落脚量（闭式律） & \cref{eq:lb-alip} \\
\texttt{numContacts} & 接触数：点足 $=2$（四足 $=4$） & \cref{tab:lb-dims} \\
\bottomrule
\end{tabular}
\end{table}

\begin{table}[H]\centering\small
\caption{公式速查表}
\begin{tabular}{lll}
\toprule
公式 / 结论 & 一句话说明 & 对应 \\
\midrule
LIP 动力学 \cref{eq:lb-lip} & $\ddot x=\omega^2(x-p)$，\emph{正号}=不稳定 & \cref{ins:lb-positive-sign} \\
DCM 定义 \cref{eq:lb-dcm-def} & $\xi=x+\dot x/\omega$，沿发散方向 & \cref{note:lb-dcm-name} \\
DCM 动力学 \cref{eq:lb-dcm-stance} & $\dot\xi=\omega(\xi-p)$，一维 & \cref{eq:lb-dcm-sol} \\
捕获点 \cref{eq:lb-capture} & $p=\xi\Rightarrow\dot\xi=0$ 一步停稳 & \cref{ins:lb-must-step} \\
Raibert 速度项 \cref{eq:lb-raibert-recall} & $\sqrt{z_0/g}(\mathbf v-\mathbf v_{des})$=捕获点反馈 & \cref{ins:lb-raibert-is-capture} \\
ALIP 横向闭式律 \cref{eq:lb-alip} & 非零步宽 period-2 极限环 & \cref{ins:lb-lateral-must} \\
DCM 投影 \cref{eq:lb-dcm-proj2} & $\exp(\mathrm{clamp}(\omega t,0,1.5))$ 外推 & \cref{pit:lb-clamp-exp} \\
机身任务拆分 \cref{eq:lb-base-split} & XY前馈+高度PD+角(旋转误差PD) & \cref{ins:lb-split-is-specialization} \\
角任务正确式 \cref{eq:lb-base-angular} & \texttt{rotationErrorInWorld}+世界系$\boldsymbol\omega$ & \cref{pit:lb-frame-bug} \\
\bottomrule
\end{tabular}
\end{table}

\begin{table}[H]\centering\small
\caption{四足栈 $\to$ 点足移植 checklist}
\begin{tabular}{p{0.46\textwidth} p{0.46\textwidth}}
\toprule
检查项 & 点足为何必须 \\
\midrule
WBC 机身任务是否带反馈（非纯前馈） & 无静稳倒立摆，纯前馈必发散（\cref{sec:lb-wbc}） \\
角任务用世界系旋转误差而非欧拉角直接相加 & 避免切空间错配 + gimbal 奇异（\cref{pit:lb-frame-bug}） \\
期望量解析得到，不做有限差分 & 防 $476\times$ 噪声放大（\cref{pit:lb-finite-diff}） \\
\texttt{frictionCoefficient} 配点足值且两层一致 & 软锥与 WBC 金字塔可行集须一致（\cref{tab:lb-braver-config}） \\
base 位姿权重（尤其 z）摁紧 & 无静稳须守站高（\cref{tab:lb-braver-config}） \\
\texttt{defaultModeSequenceTemplate} 设为踏步 & 从第一拍就走，避 $1$ 时域空窗（\cref{sec:lb-runtime}） \\
步态模板：\texttt{switchingTimes} 比模式多一个 & 否则平铺崩溃（\cref{pit:lb-tile}） \\
默认是否禁止负 $f^z$（如需） & 软锥对负法向力免费（\cref{pit:lb-soft-cone-free}） \\
\bottomrule
\end{tabular}
\end{table}

\begin{table}[H]\centering\small
\caption{知识点总表（行星号为难度）}
\begin{tabular}{clll}
\toprule
\# & 知识点 & 核心要点 & 难度 \\
\midrule
1 & 点足无静稳 & 同 $x$ 线$\Rightarrow$俯仰零力臂$\Rightarrow$倒立摆 & \(\bigstar\bigstar\) \\
2 & LIP 不稳定性 & $\ddot x=\omega^2 x$ 正号$\Rightarrow$发散模态 & \(\bigstar\bigstar\) \\
3 & DCM 分离 & $\xi=x+\dot x/\omega$，$\dot\xi=\omega\xi$ 浓缩一维 & \(\bigstar\bigstar\bigstar\) \\
4 & 捕获点 & $p=\xi$ 一步停稳；=Raibert 速度项 & \(\bigstar\bigstar\bigstar\) \\
5 & ALIP 横向闭式律 & 非零步宽 period-2；矢状步长/横向步宽 & \(\bigstar\bigstar\bigstar\bigstar\) \\
6 & OCS2 五块拼图 & ReferenceManager 的并行因果链 & \(\bigstar\bigstar\bigstar\) \\
7 & 注册$\ne$满足 & 唯一证据 \texttt{inequalityConstraintsSSE} & \(\bigstar\bigstar\bigstar\) \\
8 & 软锥两陷阱 & \texttt{mu}=惩罚权重；负 $f^z$ 免费 & \(\bigstar\bigstar\bigstar\) \\
9 & 机身任务拆分 & 角用世界系旋转误差 PD & \(\bigstar\bigstar\bigstar\bigstar\) \\
10 & 两块根因 & frame bug + 差分噪声 $476\times$ & \(\bigstar\bigstar\bigstar\bigstar\) \\
11 & 运行时步态 & 改默认序列；平铺陷阱 $N{+}1$ & \(\bigstar\bigstar\bigstar\) \\
\bottomrule
\end{tabular}
\end{table}

% ════════════════════════════════════════════════════════════════
\section*{延伸阅读}
\begin{itemize}
  \item \textbf{捕获点奠基}：Pratt 等\cite{pratt2006capture}——捕获点与可捕获区域的 push recovery 框架；落足启发式的源头见 Raibert\cite{raibert1986legged}；DCM 的三维系统化见 Englsberger 等\cite{englsberger2015dcm}。
  \item \textbf{ALIP / H-LIP 落脚律}：Gong \& Grizzle\cite{gong2022alip}（绕接触点角动量 LIP，本章横向闭式律 \cref{eq:lb-alip} 的出处）、Gibson/Gong/Grizzle\cite{gibson2022alipmpc}（落脚为决策变量的 ALIP-MPC）、Xiong \& Ames\cite{xiong2022hlip}（H-LIP）、Khadiv 等\cite{khadiv2016stepping}（步位+步时）、Acosta \& Posa\cite{acosta2023mpfc}（落脚 MPC 谱系）。
  \item \textbf{横向平衡生物力学}：Kuo\cite{kuo1999stabilization}——人靠调\emph{步宽}稳横向、靠调步长稳矢状，本章 \cref{ins:lb-lateral-must} 的依据。
  \item \textbf{点足 / 欠驱 WBC}：点足全身操作空间控制 WBOSC\cite{wbosc2015}；欠驱 point-foot biped 的 Riccati/LQR-WBC\cite{dai2024pointfoot}——本章拆分机身任务的文献对照。
  \item \textbf{质心 MPC 与简化模型谱系}：单刚体凸 MPC 见 \cref{ch:quad_mpc} 与其延伸阅读；LIPM 源头 Kajita\cite{kajita2001lipm}、质心动力学 Orin\cite{orin2013centroidal}；OCS2 感知质心 NMPC（同框架）Grandia 等\cite{grandia2022perceptive}；安全层 CBF+MPC 见 Grandia 等\cite{grandia2021cbf}；综述见 Wensing 等\cite{wensing2024legged}。
  \item \textbf{全身 MPC 对照}：基于 MuJoCo 有限差分的全身 MPC\cite{howell2023wholebodympc}——与本章“质心 MPC + WBC 分层”形成对照。
\end{itemize}

\section*{本章与后续章节的关系}
\begin{table}[H]\centering\small
\begin{tabular}{lll}
\toprule
相关章节 & 关系 & 本章接口 \\
\midrule
\cref{ch:quad_arch} 导论与控制架构 & 提供浮动基欠驱/GRF 唯一外力/多模型立论 & \cref{ins:lb-no-static} \\
\cref{ch:quad_gait} 步态与落足 & 其 Raibert 速度项即捕获点反馈 & \cref{eq:qg-raibert-vel},\cref{ins:lb-raibert-is-capture} \\
\cref{ch:quad_mpc} 单刚体 MPC & LIPM 谱系 + QP 化 + 摩擦凸化 & \cref{sec:srbd},\cref{sec:qp} \\
\cref{ch:quad_wbc} 全身控制 WBC & 浮动基子任务的点足特化 & \cref{sec:qw-floating},\cref{tab:qw-subtasks} \\
\cref{ch:lie} 李群与李代数 & 旋转误差 \texttt{rotationErrorInWorld} & \cref{eq:lb-base-angular} \\
\cref{ch:legged-debug} 调试方法论 & 本章失败根因（层级归因/规划vs执行）完整复盘 & \cref{sec:lb-wbc},\cref{pit:lb-register-ne-satisfy} \\
\bottomrule
\end{tabular}
\end{table}

\section*{故障排查手册}
\begin{enumerate}
  \item \textbf{症状}：点足总是 \emph{roll 先倒}、力矩撞饱和。\textbf{别}先怀疑摩擦/落脚/力可行性（那是\emph{后果}）。\textbf{先查}：WBC 机身任务是否带反馈（\cref{sec:lb-wbc}）；同时打印“controller 实测量 vs MPC 对未来的预测量”——若 MPC 预测收敛而实测发散（如 $50\,\mathrm{ms}$ 内 pitch 比自然倒下快约 $17$ 倍），问题 $100\%$ 在执行层，立刻停止碰上游（完整方法论见 \cref{ch:legged-debug}）。
  \item \textbf{症状}：注册了硬约束 $f^z\ge\text{floor}$，求解器报“收敛”但机器人仍不可行。\textbf{原因}：\texttt{g\_max} 不含不等式残差（\cref{pit:lb-register-ne-satisfy}）。\textbf{对策}：查 \texttt{inequalityConstraintsSSE}；若它很大，说明是参考不可行（机身失稳）逼出的，回去修机身反馈，别在不可行参考上加迭代。
  \item \textbf{症状}：调大 \texttt{frictionConeSoftConstraint.\allowbreak mu} 后求解\emph{更病态}、抖。\textbf{原因}：\texttt{mu} 是惩罚权重不是摩擦系数（\cref{pit:lb-mu-naming}）。\textbf{对策}：改摩擦红线动 \texttt{frictionCoefficient}；\texttt{mu}/\allowbreak\texttt{delta} 只调“越线罚多重”。
  \item \textbf{症状}：OCS2 desired stance 力频繁出现\emph{负} $f^z$。\textbf{原因}：默认软锥不含 $f^z\ge 0$、对负法向力免费（\cref{pit:lb-soft-cone-free}）；也可能叠加身体倒。\textbf{对策}：如需禁负，显式加 $f^z\ge 0$ 约束（注意 \cref{pit:lb-register-ne-satisfy}）。
  \item \textbf{症状}：运行时注入 gait 后机器人\emph{先愣 $\sim 1\,\mathrm{s}$ 再动}、期间倒掉。\textbf{原因}：\texttt{GaitReceiver} 在时域末端才 insert，约 $1$ 时域延迟（\cref{subsec:lb-gait-delay}）。\textbf{对策}：改 \texttt{defaultModeSequenceTemplate} 为踏步序列，从第一拍就走。
  \item \textbf{症状}：定义步态后\emph{一条腿整段 swing 不落}、或求解器越界崩溃。\textbf{原因}：单模板被平铺成连续单支撑，或 \texttt{switchingTimes} 比 \texttt{modeSequence} 少一个（\cref{pit:lb-tile}）。\textbf{对策}：模板写成 \texttt{stance$\to$DS$\to$opposite\_stance$\to$DS}，确认切换时刻数 $=$ 模式数 $+1$。
  \item \textbf{症状}：切到 \texttt{FullCentroidal} 后机身任务突然\emph{抖到发散}（SRBD 下没事）。\textbf{原因}：机身任务用了有限差分算关节加速度，被 $\sim 476\times$ 放大，SRBD 因 $\mathbf A_j\approx 0$ 幸免、Full 致命（\cref{pit:lb-finite-diff}）。\textbf{对策}：期望量改用 \texttt{computeBaseKinematicsFromCentroidalModel} 解析得到。
\end{enumerate}
